% Section 2.1, from lectures/L02/appendix/a_logic_gates.md.
%
% A.10 is headed "Kapitlets krets" here, where the lecture has "Föreläsningens krets". Two
% corrections, made in the source too: A.4 marks every law that does not hold in ordinary algebra
% (the lecture said so but left Dominans and Absorption unmarked), and A.8 no longer promises the
% multiplexer in L06, which never mentions it. A.5's laws for three variables are set as a display
% rather than inline, where they would not fit on a line.

\section{Logiska grindar och boolesk algebra}
\label{c2:app:a}

\subsection{Digitala signaler och logiska nivåer}
\label{c2:a1}
\Secref{c1:a1} konstaterade att en digital signal bara har två värden, \code{0} och \code{1}. I en
verklig krets är de två värdena två spänningsintervall:
\begin{itemize}
  \item I en krets som matas med 5~V, som IC-kretsarna i Labb~2 (bilaga~\ref{app:lab2}) och
    mikrodatorn i Labb~3, läses en spänning nära 0~V som \code{0} och en spänning nära 5~V som
    \code{1}. Exakt var gränserna går står i kretsens datablad (kapitel~\ref{c4:ch}).
  \item I ett kontaktnät på 24~V, som i Labb~1 (bilaga~\ref{app:lab1}), är en tryckknapp \code{1}
    när den är nedtryckt och \code{0} när den är släppt, och lampan är \code{1} när den lyser.
\end{itemize}

Det är en överenskommelse att hög spänning betyder \code{1}. Den kallas \textbf{positiv logik} och
är den som används i hela boken.

En \textbf{logisk grind} är en krets med en eller flera ingångar och en utgång, där utgången är en
fast funktion av ingångarna. Ett nät av grindar, där utgångarna bara beror på vad ingångarna är
just nu, kallas ett \textbf{kombinatoriskt nät}. Det är ämnet för det här kapitlet och de tre
följande. En krets som också minns vad som har hänt tidigare kallas ett \textbf{sekvensnät}; den
typen får du en glimt av i \secref{c2:b7}, och den är ämnet för kapitel~\ref{c6:ch}.

\subsection{Grindarna}
\label{c2:a2}
Det finns sju grundgrindar. NOT har en ingång; de andra har två eller fler.

\begin{booktable}{@{}p{8.6em}Lp{7.4em}@{}}
  \thead{Grind} & \thead{Utgången är \code{1} när \dots} & \thead{Uttryck} \\
  \midrule
  NOT (inverterare) & ingången är \code{0} & \code{X = A'} \\
  AND & \textbf{alla} ingångar är \code{1} & \code{X = AB} \\
  OR & \textbf{minst en} ingång är \code{1} & \code{X = A + B} \\
  NAND & inte alla ingångar är \code{1} (inverterad AND) & \code{X = (AB)'} \\
  NOR & ingen ingång är \code{1} (inverterad OR) & \code{X = (A + B)'} \\
  XOR & ingångarna är \textbf{olika} & \code{X = A ⊕ B} \\
  XNOR & ingångarna är \textbf{lika} & \code{X = (A ⊕ B)'} \\
\end{booktable}

\noindent Sanningstabellerna för alla grindar med två ingångar, sida vid sida:

\begin{narrowtable}{@{}cc@{\hspace{1.5em}}cccccc@{}}
  \thead{A} & \thead{B} & \thead{AND} & \thead{OR} & \thead{NAND} & \thead{NOR} & \thead{XOR} &
    \thead{XNOR} \\
  \midrule
  0 & 0 & 0 & 0 & 1 & 1 & 0 & 1 \\
  0 & 1 & 0 & 1 & 1 & 0 & 1 & 0 \\
  1 & 0 & 0 & 1 & 1 & 0 & 1 & 0 \\
  1 & 1 & 1 & 1 & 0 & 0 & 0 & 1 \\
\end{narrowtable}

\begin{narrowtable}{@{}cc@{}}
  \thead{A} & \thead{NOT A} \\
  \midrule
  0 & 1 \\
  1 & 0 \\
\end{narrowtable}

\noindent Värt att lägga på minnet direkt ur tabellerna:
\begin{itemize}
  \item NAND, NOR och XNOR är inverserna av AND, OR och XOR: varje kolumn är den andras med ettor
    och nollor bytta.
  \item AND och OR går att utvidga till fler ingångar. En AND med tre ingångar är \code{1} bara om
    alla tre är \code{1}; en OR med tre ingångar är \code{1} om minst en är \code{1}.
  \item XOR med fler än två ingångar är \code{1} när ett \textbf{udda} antal ingångar är \code{1}.
    Det är samma sak som paritetsbiten i \secref{c1:b6}.
\end{itemize}

\subsubsection{Två sätt att rita samma grind}
Grindarna ritas med två olika symbolstandarder, och du kommer att möta båda:
\begin{itemize}
  \item \textbf{ANSI} (amerikansk standard) har en egen form för varje grind: AND är D-formad, OR
    har en böjd ingångssida. Den används i CircuitVerse och i de flesta datablad.
  \item \textbf{IEC} (internationell standard) ritar varje grind som en rektangel med en symbol i:
    \code{&} för AND, \code{≥1} för OR (minst en ingång ska vara 1), \code{=1} för XOR (exakt en
    ingång ska vara 1) och \code{1} för NOT. Den förekommer ofta i svenska läroböcker och
    kopplingsscheman.
\end{itemize}
I båda standarderna betyder en liten ring på utgången att utsignalen är inverterad
(figur~\ref{c2:fig:gate_symbols}).

\bookfigure[tbp]{l02_gate_symbols}{De sju grundgrindarna med ANSI-symbol, IEC-symbol och booleskt
uttryck.}{c2:fig:gate_symbols}

\subsection{Boolesk algebra som notation}
\label{c2:a3}
Att skriva \enquote{A AND (NOT B) OR C} blir snabbt otympligt. \textbf{Boolesk algebra}, efter
matematikern George Boole, ger en kort notation som ser ut som vanlig algebra:

\begin{narrowtable}{@{}lll@{}}
  \thead{Operation} & \thead{Skrivs} & \thead{Utläses} \\
  \midrule
  NOT & \code{A'} & \enquote{A inverterad}, \enquote{icke A} \\
  AND & \code{AB} eller \code{A·B} & \enquote{A och B} \\
  OR & \code{A + B} & \enquote{A eller B} \\
  XOR & \code{A ⊕ B} & \enquote{A exklusivt eller B} \\
\end{narrowtable}

\noindent Du kommer att se NOT skrivet på andra sätt också, framför allt som ett streck över
bokstaven, $\overline{A}$. Boken använder primtecknet, \code{A'}, eftersom det går att skriva på ett
tangentbord.

\textbf{Prioritetsordningen} är NOT först, sedan AND, sist OR, precis som potenser, multiplikation
och addition i vanlig algebra:
\begin{itemize}
  \item \mbox{\code{AB + C}} betyder \mbox{\code{(AB) + C}}.
  \item \mbox{\code{A + BC}} betyder \mbox{\code{A + (BC)}}, och det är \textbf{inte} samma sak som
    \mbox{\code{(A + B)C}}.
  \item \code{AB'} betyder \code{A(B')}: bara B är inverterad. \code{(AB)'} inverterar hela
    produkten.
\end{itemize}
Sätt ut parenteser när du är osäker. De kostar ingenting.

Ett uttryck som är en OR av AND-termer, till exempel \mbox{\code{X = AB' + CD}}, sägs stå på
\textbf{SP-form}, summa av produkter. Det är den form nästan alla uttryck i boken skrivs på,
eftersom den läses av direkt ur en sanningstabell (\secref{c2:a6}) och översätts direkt till ett
grindnät: en AND-grind per term, och en OR-grind som samlar dem.

\subsection{Räknelagarna}
\label{c2:a4}
Lagarna nedan gäller för alla värden på A, B och C, och varje lag går att kontrollera genom att
skriva ut sanningstabellen för båda sidorna. Många av dem ser ut som vanlig algebra; några gör det
inte, och de är markerade.

\begin{booktable}{@{}>{\raggedright\arraybackslash}p{8.4em}LL@{}}
  \thead{Lag} & \thead{Med OR} & \thead{Med AND} \\
  \midrule
  Identitet & \code{A + 0 = A} & \code{A · 1 = A} \\
  Dominans & \code{A + 1 = 1} \emph{(inte som vanlig algebra)} & \code{A · 0 = 0} \\
  Idempotens \emph{(inte som vanlig algebra)} & \code{A + A = A} & \code{A · A = A} \\
  Komplement & \code{A + A' = 1} & \code{A · A' = 0} \\
  Dubbel invertering & \code{(A')' = A} & \\
  Kommutativ & \code{A + B = B + A} & \code{AB = BA} \\
  Associativ & \code{A + (B + C) = (A + B) + C} & \code{A(BC) = (AB)C} \\
  Distributiv & \code{A(B + C) = AB + AC} & \code{A + BC = (A + B)(A + C)} \emph{(inte som vanlig
    algebra)} \\
  Absorption \emph{(inte som vanlig algebra)} & \code{A + AB = A} & \code{A(A + B) = A} \\
  Förenkling & \code{A + A'B = A + B} & \code{A(A' + B) = AB} \\
\end{booktable}

\noindent Några av lagarna blir självklara när man tänker på vad de betyder:
\begin{itemize}
  \item \mbox{\code{A + 1 = 1}}: \enquote{A eller sant} är alltid sant.
  \item \mbox{\code{A + A' = 1}}: en signal är alltid antingen \code{1} eller \code{0}, så
    \enquote{A eller icke A} är alltid sant.
  \item \mbox{\code{A + AB = A}}: om A är \code{1} är uttrycket \code{1} oavsett B, och om A är
    \code{0} är båda termerna \code{0}. Termen \code{AB} tillför alltså ingenting.
\end{itemize}

Den sista raden, \textbf{förenklingslagen} \mbox{\code{A + A'B = A + B}}, är den som oftast behövs
och oftast missas. Den går att härleda ur den andra distributiva lagen:

\begin{textdrawing}
A + A'B = (A + A')(A + B) = 1 · (A + B) = A + B
\end{textdrawing}

\noindent Med ord: i termen \code{A'B} behövs inte villkoret \code{A'}, för när A är \code{1} är
uttrycket redan \code{1}.

\subsubsection{Att bevisa en lag med en sanningstabell}
Varje lag kan kontrolleras genom att räkna ut båda sidorna för alla kombinationer.
Förenklingslagen:

\begin{narrowtable}{@{}cccccc@{}}
  \thead{A} & \thead{B} & \thead{A'} & \thead{A'B} & \thead{A + A'B} & \thead{A + B} \\
  \midrule
  0 & 0 & 1 & 0 & 0 & 0 \\
  0 & 1 & 1 & 1 & 1 & 1 \\
  1 & 0 & 0 & 0 & 1 & 1 \\
  1 & 1 & 0 & 0 & 1 & 1 \\
\end{narrowtable}

\noindent De två sista kolumnerna är lika på varje rad, så lagen gäller. Med två variabler är det
fyra rader och med tre är det åtta, så det finns ingen ursäkt för att inte kontrollera en förenkling
man är osäker på.

\subsection{De Morgans lagar}
\label{c2:a5}
Två lagar är så viktiga att de har ett eget namn, efter matematikern Augustus De Morgan:

\begin{console}
(A + B)' = A'B'
(AB)'    = A' + B'
\end{console}

\noindent Med ord: \textbf{invertera hela uttrycket genom att invertera varje variabel och byta AND
mot OR} (och tvärtom). En minnesregel är \enquote{bryt strecket, byt tecknet}.

Kontrollera den första med en tabell:

\begin{narrowtable}{@{}ccccccc@{}}
  \thead{A} & \thead{B} & \thead{A + B} & \thead{(A + B)'} & \thead{A'} & \thead{B'} &
    \thead{A'B'} \\
  \midrule
  0 & 0 & 0 & 1 & 1 & 1 & 1 \\
  0 & 1 & 1 & 0 & 1 & 0 & 0 \\
  1 & 0 & 1 & 0 & 0 & 1 & 0 \\
  1 & 1 & 1 & 0 & 0 & 0 & 0 \\
\end{narrowtable}

\noindent Kolumnerna \mbox{\code{(A + B)'}} och \code{A'B'} är lika. Med ord:
\enquote{inte (A eller B)} är samma sak som \enquote{varken A eller B}, alltså
\enquote{inte A och inte B}.

Lagarna gäller för fler variabler också:

\begin{console}
(A + B + C)' = A'B'C'
(ABC)'       = A' + B' + C'
\end{console}

\subsubsection{Vad De Morgan betyder för grindarna}
De Morgan är mer än en räkneregel. Den säger att en grind kan bytas mot en annan grind med
inverterade in- och utgångar:
\begin{itemize}
  \item En \textbf{NAND}, \code{(AB)'}, är samma sak som en \textbf{OR med båda ingångarna
    inverterade}, \mbox{\code{A' + B'}}.
  \item En \textbf{NOR}, \mbox{\code{(A + B)'}}, är samma sak som en \textbf{AND med båda ingångarna
    inverterade}, \code{A'B'}.
\end{itemize}

Det har två viktiga följder. Den första är att \textbf{varje funktion går att bygga med enbart
NAND-grindar}, eller med enbart NOR-grindar. En NAND med båda ingångarna ihopkopplade är en NOT,
eftersom \mbox{\code{(AA)' = A'}}; en NAND följd av en NOT är en AND; och en OR är en NAND med
inverterade ingångar. Det är därför NAND-kretsen är den vanligaste IC-kretsen av alla, och det är
vad Labb~2-2 (\secref{lab2:2}) går ut på.

Den andra följden syns i kontaktnät: en NAND är två brytande kontakter parallellt, och en NOR två
brytande kontakter i serie (\secref{c2:b5}).

\subsection{Från sanningstabell till uttryck}
\label{c2:a6}
Ofta börjar ett konstruktionsarbete med en beskrivning av vad kretsen ska göra, och den
beskrivningen blir en sanningstabell. Varje sanningstabell kan läsas av direkt som ett uttryck på
SP-form:
\begin{enumerate}
  \item Leta upp varje rad där utgången är \code{1}.
  \item Skriv för varje sådan rad en AND-term med \textbf{alla} ingångar: ingången som den är om
    den är \code{1} på raden, inverterad om den är \code{0}.
  \item Sätt ihop termerna med OR.
\end{enumerate}
En sådan AND-term, med alla ingångar och som är \code{1} på exakt en rad, kallas en
\textbf{minterm}.

\textbf{Exempel:}

\begin{narrowtable}{@{}ccc@{}}
  \thead{A} & \thead{B} & \thead{X} \\
  \midrule
  0 & 0 & 0 \\
  0 & 1 & 1 \\
  1 & 0 & 1 \\
  1 & 1 & 0 \\
\end{narrowtable}

\noindent Raden \mbox{\code{AB = 01}} ger mintermen \code{A'B}, och raden \mbox{\code{AB = 10}} ger
\code{AB'}. Alltså:

\begin{console}
X = A'B + AB'
\end{console}

\noindent Det är XOR-funktionen. Samma nät går alltså att bygga av en enda XOR-grind, i stället för
två AND, två NOT och en OR.

\textbf{Exempel med tre ingångar}, en \textbf{majoritetsgrind}: utgången är \code{1} när minst två
av de tre ingångarna är \code{1}. Den används där tre givare mäter samma sak och systemet ska lita
på det två av dem är överens om, så att en enstaka trasig givare inte kan bestämma något på egen
hand.

\begin{narrowtable}{@{}ccccl@{}}
  \thead{A} & \thead{B} & \thead{C} & \thead{X} & \thead{Minterm} \\
  \midrule
  0 & 0 & 0 & 0 & \\
  0 & 0 & 1 & 0 & \\
  0 & 1 & 0 & 0 & \\
  0 & 1 & 1 & 1 & \code{A'BC} \\
  1 & 0 & 0 & 0 & \\
  1 & 0 & 1 & 1 & \code{AB'C} \\
  1 & 1 & 0 & 1 & \code{ABC'} \\
  1 & 1 & 1 & 1 & \code{ABC} \\
\end{narrowtable}

\begin{console}
X = A'BC + AB'C + ABC' + ABC
\end{console}

\noindent Uttrycket är \textbf{korrekt}, men det är inte det \textbf{minsta}: fyra AND-grindar med
tre ingångar och en OR-grind med fyra. Att förenkla det är nästa avsnitts ämne.

\subsection{Algebraisk förenkling}
\label{c2:a7}
Ett uttryck som läses av rakt ur en sanningstabell har en term per etta. Ofta delar flera termer
det mesta av sina variabler, och då går de att slå ihop. Verktyget är nästan alltid samma:
\textbf{bryt ut det gemensamma, och använd \mbox{\code{A + A' = 1}}}.

\textbf{Exempel 1:}

\begin{textdrawing}
X = AB + AB' = A(B + B') = A · 1 = A
\end{textdrawing}

\noindent Värdet på B spelar ingen roll, så B försvinner.

\textbf{Exempel 2:} en sanningstabell där X är \code{1} på raderna 001, 011, 101 och 111.

\begin{textdrawing}
X = A'B'C + A'BC + AB'C + ABC
  = A'C(B' + B) + AC(B' + B)
  = A'C + AC
  = C(A' + A)
  = C
\end{textdrawing}

\noindent Fyra termer blev en enda variabel. Titta på tabellen igen: X är \code{1} precis när C är
\code{1}. Det hade gått att se direkt, men uträkningen visar varför.

\textbf{Exempel 3:} majoritetsgrinden från \secref{c2:a6}. Här behövs ett knep: termen \code{ABC}
kan användas flera gånger, eftersom \mbox{\code{ABC + ABC = ABC}} (idempotens). Skriv den tre
gånger, och para ihop den med var och en av de andra termerna:

\begin{textdrawing}
X = A'BC + AB'C + ABC' + ABC
  = (A'BC + ABC) + (AB'C + ABC) + (ABC' + ABC)
  = BC(A' + A)  + AC(B' + B)  + AB(C' + C)
  = BC + AC + AB
\end{textdrawing}

\noindent Tre AND-grindar med två ingångar och en OR-grind med tre, i stället för fyra AND-grindar
med tre ingångar.

\textbf{Exempel 4:} inte allt går att förenkla. \mbox{\code{X = A'B + AB'}} har inga gemensamma
variabler att bryta ut. Uttrycket är redan minimalt på SP-form, och det bästa man kan göra är att
känna igen det som en XOR.

Algebraisk förenkling kräver att man ser vilka termer som går att para ihop, och det är lätt att
missa ett par, eller att inte veta när man är klar. Kapitel~\ref{c4:ch} introducerar
\textbf{Karnaughdiagram}, som gör samma sak med en ritning, systematiskt, och med ett tydligt svar
på när man är klar.

\subsection{Analys: från grindnät till sanningstabell}
\label{c2:a8}
Åt andra hållet: du har ett färdigt grindnät och ska ta reda på vad det gör. Metoden är att ge
varje grinds utgång ett namn, en \textbf{mellansignal}, och räkna fram dem en i taget, rad för rad.

\bookfigure{l02_analysis_network}{Ett grindnät att analysera, med mellansignalerna A', P och
Q.}{c2:fig:analysis_network}

\noindent Mellansignalerna är \code{A'}, \mbox{\code{P = AB}} och \mbox{\code{Q = A'C}}, och
utgången är \mbox{\code{X = P + Q}}. Uttrycket för hela nätet blir alltså

\begin{console}
X = AB + A'C
\end{console}

\noindent och sanningstabellen räknas fram kolumn för kolumn:

\begin{narrowtable}{@{}ccccccc@{}}
  \thead{A} & \thead{B} & \thead{C} & \thead{A'} & \thead{P = AB} & \thead{Q = A'C} &
    \thead{X = P + Q} \\
  \midrule
  0 & 0 & 0 & 1 & 0 & 0 & 0 \\
  0 & 0 & 1 & 1 & 0 & 1 & 1 \\
  0 & 1 & 0 & 1 & 0 & 0 & 0 \\
  0 & 1 & 1 & 1 & 0 & 1 & 1 \\
  1 & 0 & 0 & 0 & 0 & 0 & 0 \\
  1 & 0 & 1 & 0 & 0 & 0 & 0 \\
  1 & 1 & 0 & 0 & 1 & 0 & 1 \\
  1 & 1 & 1 & 0 & 1 & 0 & 1 \\
\end{narrowtable}

\noindent Läs tabellen i två halvor. När A är \code{0} är X samma som C; när A är \code{1} är X
samma som B. Nätet är alltså en \textbf{väljare}: A väljer vilken av B och C som ska släppas fram
till utgången. En sådan krets kallas en \textbf{multiplexer}, och den är ett av de viktigaste
byggblocken i en dator. Den återkommer i övning~\ref{c4:ex:13}.

Två råd som gäller för all analys:
\begin{itemize}
  \item \textbf{Skriv ut mellansignalerna som egna kolumner.} Att räkna fram X direkt ur A, B och
    C i huvudet fungerar för små nät och ger fel i större.
  \item \textbf{Sammanfatta sedan tabellen med ord.} En tabell säger vad kretsen gör rad för rad;
    en mening som \enquote{A väljer mellan B och C} säger vad den är till för.
\end{itemize}

\subsection{CircuitVerse}
\label{c2:a9}
CircuitVerse, \url{https://circuitverse.org/simulator}, är en gratis logiksimulator som körs i
webbläsaren. Varje grindnät i boken byggs och testas där innan det byggs med riktiga kretsar:
\begin{enumerate}
  \item Öppna simulatorn och starta ett nytt projekt.
  \item Placera ett \textbf{Input}-element per insignal och ett \textbf{Output}-element per
    utsignal, och döp om dem så att de heter som dina variabler.
  \item Placera de grindar ditt uttryck kräver, från panelen \textbf{Gates}.
  \item Koppla ihop nätet en grind i taget, enligt uttrycket. Du ska redan veta exakt vilka grindar
    du behöver innan du öppnar simulatorn, och det är därför \secref{c2:a6} och \ref{c2:a7} kommer
    först.
  \item Klicka på varje ingång för att växla den mellan \code{0} och \code{1}, och kontrollera att
    utgången stämmer med din sanningstabell på \textbf{varje} rad.
  \item Spara under menyn \textbf{Project}, antingen online med ett konto eller som en fil på
    datorn.
\end{enumerate}

\textbf{Steg 5 är inte valfritt, och ingenting annat gör det åt dig.} Tre regler gör det till en
kontroll snarare än en formalitet:
\begin{itemize}
  \item \textbf{Alla rader, inte ett urval.} Två ingångar är fyra rader och tre är åtta. Vid de
    storlekarna finns ingen ursäkt för att bara prova några.
  \item \textbf{Förutsäg först, simulera sedan.} Skriv sanningstabellen innan du bygger. Annars
    läser du av simuleringen och tror att du har kontrollerat den, när du bara har skrivit av den.
  \item \textbf{En avvikelse är information.} Den säger att ritningen och uttrycket är oense, och
    att en av dem har ett fel.
\end{itemize}

Den vanan, att förutsäga ett beteende och sedan kontrollera det, är den du kommer att använda vid
varje laboration i boken, med kontakter, med IC-kretsar och med mikrodatorn.

\subsection{Kapitlets krets: bromsassistenten}
\label{c2:a10}
Kapitlet bygger en krets från ett krav till ett simulerat grindnät, och i \secref{c2:b6} samma krets
som kontaktnät.

Kravet: \textbf{bilen bromsar om föraren bromsar, eller om assistanssystemet bestämmer sig för
det.} Assistanssystemet bestämmer sig för det när endera av två detektorer ser ett hinder, om det
inte samtidigt rapporterar ett fel, i vilket fall det inte anförtros att bromsa alls.

\begin{booktable}{@{}p{5em}L@{}}
  \thead{Signal} & \thead{Betydelse} \\
  \midrule
  \code{D} & Föraren står på bromspedalen. \\
  \code{S} & Avståndssensorn ser ett hinder. \\
  \code{R} & Radarn ser ett hinder. \\
  \code{F} & Assistanssystemet rapporterar ett fel. \\
  \code{B} & Utgång: lägg an bromsarna. \\
\end{booktable}

\noindent Tre meningar i kravet bestämmer tillsammans hela strukturen:
\begin{itemize}
  \item Endera detektorn räcker på egen hand, och ingen av dem väger tyngre än den andra.
  \item Ett fel utlöser inget larm; det \emph{hindrar} assistanssystemet från att bromsa.
  \item Föraren kan alltid bromsa, även när assistanssystemet är trasigt.
\end{itemize}
Den sista är ett säkerhetskrav, och det är den att hålla fast vid: pedalen är inte en ingång till
assistanslogiken, den går förbi den.

\subsubsection{Sanningstabellen}
Sexton rader, i den ordning en fyrabitars räknare skulle producera dem. Kolumnen \code{P} är ingen
utgång; den är assistanssystemets eget beslut, visad för att det ska synas var strukturen kommer
ifrån.

\begin{narrowtable}{@{}cccc@{\hspace{1.5em}}cc@{}}
  \thead{D} & \thead{S} & \thead{R} & \thead{F} & \thead{P} & \thead{B} \\
  \midrule
  0 & 0 & 0 & 0 & 0 & 0 \\
  0 & 0 & 0 & 1 & 0 & 0 \\
  0 & 0 & 1 & 0 & 1 & 1 \\
  0 & 0 & 1 & 1 & 0 & 0 \\
  0 & 1 & 0 & 0 & 1 & 1 \\
  0 & 1 & 0 & 1 & 0 & 0 \\
  0 & 1 & 1 & 0 & 1 & 1 \\
  0 & 1 & 1 & 1 & 0 & 0 \\
  1 & 0 & 0 & 0 & 0 & 1 \\
  1 & 0 & 0 & 1 & 0 & 1 \\
  1 & 0 & 1 & 0 & 1 & 1 \\
  1 & 0 & 1 & 1 & 0 & 1 \\
  1 & 1 & 0 & 0 & 1 & 1 \\
  1 & 1 & 0 & 1 & 0 & 1 \\
  1 & 1 & 1 & 0 & 1 & 1 \\
  1 & 1 & 1 & 1 & 0 & 1 \\
\end{narrowtable}

\noindent Den nedre halvan är genomgående \code{1}: så snart D är \code{1} ändrar ingenting annat
svaret. Det är hur \enquote{åsidosätter allt} ser ut i en tabell.

\begin{aside}
\textbf{Försök själv först.} Innan du läser vidare: läs av uttrycket på SP-form, förenkla det, och
räkna grindarna. Det raka uttrycket har elva termer; det förenklade får plats på en rad.
\end{aside}

\subsubsection{Härledningen}
Att läsa av tabellen rakt av ger elva mintermer, en per etta, var och en med fyra variabler.
Korrekt, men ohanterligt. Dela i stället upp tabellen efter D.

\textbf{Den nedre halvan (D = 1)} är \code{1} på alla åtta raderna. De åtta mintermerna tillsammans
täcker alla kombinationer av S, R och F, så de förenklas till bara \code{D}, precis som i exempel 2
i \secref{c2:a7}.

\textbf{Den övre halvan (D = 0)} är \code{1} på tre rader: SRF = 010, 100 och 110, alltså när F är
\code{0} och minst en av S och R är \code{1}:

\begin{console}
D'S'RF' + D'SR'F' + D'SRF' = D'F'(S'R + SR' + SR) = D'F'(S + R)
\end{console}

\noindent (\mbox{\code{S'R + SR' + SR}} är \code{1} när minst en av S och R är \code{1}, alltså
\mbox{\code{S + R}}; övning~\ref{c2:ex:6} c) låter dig visa det algebraiskt.)

Tillsammans:

\begin{console}
B = D + D'(S + R)F'
\end{console}

\noindent och förenklingslagen \mbox{\code{A + A'X = A + X}} från \secref{c2:a4} tar bort \code{D'}:

\begin{console}
B = D + (S + R)F'
\end{console}

\noindent Uttrycket läses som kravet: bromsa om föraren bromsar, eller om (sensorn eller radarn)
ser något och det inte finns något fel.

\subsubsection{Grindnätet}
Fyra grindar: en OR för \mbox{\code{S + R}}, en NOT för \code{F'}, en AND för assistanssystemets
beslut, och en OR som lägger till förarens pedal.

\bookfigure{l02_adas_network}{Bromsassistentens grindnät.}{c2:fig:adas_network}

\textbf{Säkerhetsargumentet vilar på den sista OR-grinden.} Förarens pedal går direkt in där, förbi
allt som har med felet att göra. Ett nät som i stället matar in pedalen före fellogiken,
\mbox{\code{B = (D + S + R)F'}}, är en grind enklare och ser rimligt ut, men det bromsar inte när
föraren trampar på pedalen samtidigt som assistanssystemet rapporterar fel. Övning~\ref{c2:ex:13}
låter dig hitta exakt vilka rader som går fel.

Bygg nätet i CircuitVerse och kontrollera det mot alla sexton raderna i tabellen.
